\documentclass[11pt]{article}
\usepackage{fontspec}
\setmainfont{Times New Roman}
\setsansfont{Arial}
\setmonofont{Consolas}
\usepackage{amsmath,amssymb}
\usepackage{geometry}
\usepackage{microtype}
\geometry{a4paper,margin=3cm}
\setlength{\parindent}{1.25em}
\setlength{\parskip}{0pt}
\emergencystretch=4em
\allowdisplaybreaks
\let\A\relax
\newcommand{\A}{\mathfrak A}
\newcommand{\R}{\mathfrak R}

\begin{document}

\begin{center}
\textbf{\S\ 23. \quad Formy 10. poręda (sistem $s_{III}$).}
\end{center}

Svijanje $s$ s formami 6. poręda (sistem III) (\S\ 11).

\noindent\textit{Nezvodime formy:}
\[
\begin{array}{r@{\quad}c}
\text{A)}&
\begin{array}{c|c}
(\vartheta^2\nu x)^3s_{\nu}&\cdot\\[0.45em]
\cdot&(\vartheta^2\nu x)^3s_{\nu}s_\vartheta,\ \theta_{\nu}(\vartheta^2\nu x)^3s_\vartheta
\end{array}
\end{array}
\]
\[
\begin{array}{r@{\quad}c}
\text{B)}&
\begin{array}{c|ccc}
\cdot&(a_{\nu}(j\nu x),s,u)^2&(a_{\nu}(j\nu x),s,u)^3&(a_{\nu}(j\nu x),s,u)^4\\[0.45em]
a_{\nu}(j\nu x)s_{\nu}&\cdot&(a_{\nu}^2(j\nu x),s,u)^2&(a_{\nu}^2(j\nu x),s,u)^3
\end{array}
\end{array}
\]
\[
\begin{array}{r@{\quad}c}
\text{C)}&
\begin{array}{c|c|c}
\cdot&\cdot&((\theta Hu),s,u)^2,\ ((\theta Hu)^2,s,u)\\[0.35em]
\cdot&((\vartheta\eta x),s,u)^2,\ (\theta Hu)s_\eta,\ (\theta_\eta su)&((\vartheta\eta x),s,u)^3,\ (\theta_\eta su)^2\\[0.35em]
((\vartheta\eta x)^2,s,u)&\cdot&(\theta_\eta^2su)\\[0.35em]
\cdot&(\theta_\eta(\vartheta\eta x)^2,s,u)&\cdot
\end{array}
\end{array}
\]
\[
\begin{array}{c|c}
((\theta Hu),s,u)^3,\ ((\theta Hu)^2,s,u)^2&((\theta Hu),s,u)^4\\[0.45em]
(H_\vartheta(\theta Hu),s,u)^2,\ (\theta_\eta(\theta Hu),s,u)^2,\ (\theta_\eta(\theta Hu)^2,s,u)&(\theta_\eta su)^4,\ (\theta_\eta(\theta Hu),s,u)^3\\[0.45em]
\cdot&(\theta_\eta^2(\theta Hu),s,u)^2
\end{array}
\]

\noindent\textit{Redukcije:}

A) formy $((\vartheta^2\nu x)^3,s,u)^2$, $((\vartheta^2\nu x)^3,s,u)^3$ razkladajut se:
\[
(\vartheta\nu x)(\widehat{\vartheta\nu su})(\widehat{\vartheta'\nu su})
=(\vartheta\nu x)s_\vartheta s_{\vartheta'}-(\vartheta\nu x)s_{\nu}s_\vartheta
=-2\theta(\vartheta\nu x)s_{\nu}s_\vartheta+(\vartheta\nu x)s_\vartheta^2.
\]
Ostale formy: abo imajut reducenty kako faktory, abo sut jednokratne svijanja s razkladnymi formami.

B) Reducent $a_{\nu}^2s_{\nu}$ i $(\widehat{j\nu su})s_{\nu}$.

C) 1. Ręd form $(\theta Hu)^2s_\eta$: a) reducent $(aHu)^2(asu)s_\eta$; b) dvojna redukcija rędov form $(\theta gu)^3g_{\nu}$ i $g_{\nu}(agu)^3(bgu)^2(abu)$.

Iz togo redukcija vyših form $(\theta_\eta su)^3$, $(H_\vartheta(\theta Hu),s,u)^3$ i $(a_{\nu}\cdot b_{\nu_1}^2,s,u)^4$ [poslědnja forma vměsto $(H_\vartheta su)^4$].

2. $((\vartheta\eta x),s,u)^4$: reducent $(a_\eta su)^4$.

3. $((\vartheta\eta x)^2,s,u)^3$: reducent $s_\vartheta^2s_{\nu}$ iz $(\widehat{\vartheta\eta}su)^2(Hsu)$. Formy $(\vartheta\eta x)s_\eta$, $(\vartheta\eta x)^2s_\eta$, $((\vartheta\eta x)^2,s,u)^2$, $\theta_\eta s_\eta$ razkladajut se črez svijanje s razkladnoj formoj $a_\eta s_\eta$.

4. Rędy form $H_\vartheta(\theta Hu)s_\eta$, $\theta_\eta(\theta Hu)s_\eta$, $H_\vartheta(\vartheta\eta x)s_\eta$, $\theta_\eta(\vartheta\eta x)s_\eta$: reducenty $a_\eta(aHu)s_\eta$ i $a_\eta^2s_\eta$.

Ręd form $(\theta_\eta(\vartheta\eta x),s,u)^2$: reducent $a_\eta^2(Hsu)^2$ [s izključenjem $(\theta_\eta^2(\theta Hu),s,u)^2$, ktore nastavaje iz člena $a_\eta^2(bsu)(Hsu)$ črez svijanje].

5. $(H_\vartheta(\vartheta\eta x),s,u)$, $(H_\vartheta(\vartheta\eta x),s,u)^2$: dvojna redukcija $a_\eta(aHu)s_\eta(abu)$ i $g_\vartheta(\theta gu)g_{\nu}$.

Ręd form $(H_\vartheta(\vartheta\eta x),s,u)^3$: reducent $((\vartheta\nu x)^2,s,u)^2s_{\nu}$.

6. $(H_\vartheta(\theta Hu),s,u)$ razklada se črez jednokratno svijanje s razkladnoj formoj $(\theta_{\nu}(\vartheta\nu x),s,u)$ po \S\ 3. 2), c); t. j. črez dvojnu redukciju nižšej razkladnoj formy $(H_\vartheta su)^2$: \footnote{Znak $\equiv$ znači, že produkty form nižšego stepena sut izpuščene.}
\[
\begin{aligned}
0&\equiv \theta_{\nu}(\vartheta\nu\nu_1)(\theta su)+\frac23\theta_{\nu}^2(\vartheta\nu_1x)(\theta su)+\theta_{\nu}(\vartheta\nu x)(\theta su)s_{\nu_1}\\
&\equiv -\frac12H_\vartheta(\theta Hu)(\theta su)+\theta_\eta(\theta su)(Hsu)+\frac12H_\vartheta(\theta su)(Hsu)\\
&\equiv -\frac12H_\vartheta(\theta Hu)(\theta su)+\theta_\eta(\theta su)(Hsu)+\frac12[H_\vartheta]_{\widehat{su}^2}-\frac12\cdot\frac35H_\vartheta(\theta Hu)(\theta su)\\
&\equiv -\frac35H_\vartheta(\theta Hu)(\theta su).
\end{aligned}
\]
Črez dvojnu redukciju po 1. i 5. nastavajut formuly:
\enlargethispage{\baselineskip}
\begin{align*}
0&=\theta_\eta(\theta su)(Hsu)^2+\frac14H_\vartheta(\theta Hu)(\theta su)^2-\frac34\theta_\eta(\theta Hu)(\theta su)(Hsu)\\
&\quad-\frac54\theta_\eta(\theta Hu)^2(\theta su)-(asu)^3\cdot b_\eta(bHu)^2-a_{\nu}(asu)^3\cdot b_{\nu_1}^2,\\
0&=H_\vartheta(\theta Hu)(\theta su)^3-7\theta_\eta(\theta Hu)(\theta su)^2(Hsu),\tag{31}\\
0&=a_{\nu}(asu)^2b_{\nu_1}^{2}(bsu)^2-\theta_\eta(\theta su)^2(Hsu)^2+6\theta_\eta(\theta Hu)(\theta su)^2(Hsu),\\
0&\equiv (H_\vartheta(\vartheta\eta x),s,u),\qquad
0\equiv H_\vartheta(\widehat{\vartheta\eta su})(Hsu)-4\theta_\eta^2(Hsu).
\end{align*}

\noindent\textit{Poslědki:}

1) $(\theta Hu)^2s_\eta$, $H_\vartheta(\theta Hu)s_\eta$, $\theta_\eta(\theta Hu)s_\eta$, $H_\vartheta(\vartheta\eta x)s_\eta$, $\theta_\eta(\vartheta\eta x)s_\eta$, $(H_\vartheta(\vartheta\eta x),s,u)$, $(\theta_\eta(\vartheta\eta x),s,u)^2$ sut reducenty.

2) $s^2$ ne vhodi v svijanje s formami $(5,\lambda)$ i t. d. 6. poręda.

\end{document}
